\begin{examplebox}
\textbf{例题 1（直接求最大焦角）。}
已知椭圆 $C:\frac{x^2}{4}+\frac{y^2}{3}=1$，
求 $\angle F_1PF_2$ 的最大值。

\textit{解：} $a=2$、$b=\sqrt3$、$c=1$。
由 $|y_0|\leq\sqrt3$ 及任意点半角公式，
\[
\tan\frac\theta2=\frac{|y_0|}{3}\leq\frac1{\sqrt3}.
\]
当且仅当 $P=(0,\pm\sqrt3)$ 时取等号，故
$\theta_{\max}=2\arctan(1/\sqrt3)=60^\circ$。
\textbf{余弦定理核验：}短轴端点处 $r_1=r_2=2$，所以
$\cos\theta_{\max}=(4+4-4)/(2\cdot2\cdot2)=1/2$，
同样得到 $60^\circ$。

\medskip
\textbf{例题 2（直角焦角存在的离心率条件）。}
若椭圆 $\frac{x^2}{a^2}+\frac{y^2}{b^2}=1$ 上存在
$P$ 使 $\angle F_1PF_2=90^\circ$，求离心率范围。

\textit{解：}焦角从长轴端点的 $0$ 连续变化到短轴端点的
$\theta_{\max}$，因此存在直角焦角当且仅当
\[
 \theta_{\max}\geq90^\circ
 \iff\tan\frac{\theta_{\max}}2=\frac cb\geq1
 \iff c\geq b.
\]
由 $b^2=a^2-c^2$，得到 $2c^2\geq a^2$，即
\[
 \boxed{\frac{\sqrt2}{2}\leq e<1}.
\]
等号时直角焦角恰在短轴端点；严格大于时，椭圆上有
$|y_0|=b^2/c<b$ 的对应点。

\medskip
\textbf{例题 3（已知焦角反求点）。}
椭圆 $\frac{x^2}{9}+\frac{y^2}{4}=1$ 上，
求满足 $\angle F_1PF_2=90^\circ$ 的所有点 $P$。

\textit{解：} $a=3$、$b=2$、$c=\sqrt5$。
由 $\tan45^\circ=1$，
\[
  \frac{\sqrt5\,|y_0|}{4}=1
  \Longrightarrow |y_0|=\frac4{\sqrt5}.
\]
代入椭圆方程，
\[
  x_0^2=9\left(1-\frac{(4/\sqrt5)^2}{4}\right)
       =\frac95.
\]
故共有四个点：
\[
 \boxed{P\left(\pm\frac3{\sqrt5},\,\pm\frac4{\sqrt5}\right)}
\]
其中两个正负号\textbf{独立选取}。

\medskip
\textbf{例题 4（辨认错误公式的反例）。}
椭圆 $\frac{x^2}{5}+\frac{y^2}{4}=1$ 上取
$P(\frac{\sqrt{15}}2,1)$，此时 $c=1$、$b^2=4$。
正确公式给出
\[
 \boxed{\tan\frac\theta2=\frac{1\cdot1}{4}=\frac14}.
\]
若误用 $c/|y_0|$，会得到 $\tan(\theta/2)=1$，即
$\theta=90^\circ$；但此椭圆的最大焦角只有
$2\arctan(1/2)<90^\circ$，直接矛盾。
因此 $c/|y_0|$ 不能推广到任意点。
\end{examplebox}
